
\section{Validation: is the cleaning gain real, and do the flags hold up?}
\label{sec:validation}
The audit's value rests on two questions a skeptic should ask: is the
cleaning gain caused by better labels rather than an easier dataset, and
would the same flags appear under a different run? Three experiments
answer both, plus a calibration check on the deployed scores. All use the
untouched official 624-film test set and committed probability dumps.

\subsection{Calibration of the deployed probabilities}
On the official test set, the tuned CNN scores ECE 0.1253 /
Brier 0.1339 (skill 0.43
vs the base-rate prior) and the tuned GCN ECE 0.1413 /
Brier 0.1215 (skill 0.48): usable but only
moderately calibrated, and we say so. The zero-shot adult-domain transfer
model scores ECE 0.6126 with \textbf{negative} Brier skill
(-1.58) --- on these paediatric films its
probabilities are worse than predicting the base rate, a quantitative
statement of the domain-shift finding that threshold accuracy alone
understates (\texttt{results/calibration.json}; two independent
implementations agree to $10^{-9}$).

\subsection{Random-removal ablation: is the cleaning gain attributable to labels?}
Removing 155 \emph{random} training films under the identical split,
training configuration and test set (two seeds) gives
CNN 83.25$\pm$3.74\%
and GCN 82.77$\pm$1.93\%;
removing the \emph{audit-flagged} 155 gives CNN 79.81\% and
GCN 84.94\% (\texttt{results/pneumonia/removal\_ablation.json},
baseline 72.12\%/79.01\%). The audit's choices do \textbf{not} beat random removal: same budget, same protocol, same test, and random removal lands within one spread of flagged removal (or above it). The cleaning gain is therefore \textbf{unattributable to label quality} at $n{=}155/4{,}710$ - small-sample removal helps regardless of which films leave. We pre-registered this demotion criterion in the discussion chapter and it fires: the +5.9-point cleaning gain is demoted from finding to anecdote, and the audit's evidentiary weight rests on the gated census, its directionality, and its cross-architecture behaviour - not on a downstream accuracy claim. This negative is reported in the abstract and every section that previously asserted a causal cleaning mechanism.
\begin{table}[h]
\centering\small
\caption{Removal ablation on the official 624-film test set.}
\label{tab:removal}
\begin{tabular}{l r r r r}
\hline
removed set & CNN acc & CNN AUC & GCN acc & GCN AUC \\ \hline
random (seed 7) & 85.90\% & 93.8286 & 84.13\% & 92.3603 \\
random (seed 123) & 80.61\% & 93.8308 & 81.41\% & 93.0221 \\
\hline
\textbf{audit-flagged (155)} & \textbf{79.81\%} & \textbf{92.7833} & \textbf{84.94\%} & \textbf{93.7673} \\
\hline
\end{tabular}
\end{table}

\subsection{Seed-stability of the candidate flags}
Re-running the census at a second seed (same 3-fold assignment logic
and admissibility gate; reduced 4-epoch OOF protocol vs the canonical
8-epoch run, honestly recorded) yields the flag set in
Table~\ref{tab:stab}. The 91 images flagged in both runs
(canonical seed 42 and rerun seed 7) form a high-confidence consensus core; flags outside the core
are published as lower-confidence candidates, tiered in the artifact ---
the honest reading is that the noise \emph{rate} is seed-stable while
marginal flag \emph{identity} is representation-dependent, exactly the
structure the malaria cross-architecture study found independently.
\begin{table}[h]
\centering\small
\caption{Census seed-stability vs the canonical seed-42 census (155 flags).}
\label{tab:stab}
\begin{tabular}{r r r r r}
\hline
seed & OOF acc & flags & overlap w.\ 155 & Jaccard \\ \hline
7 & 87.18\% & 221 & 91 & 0.319 \\
\hline
\end{tabular}
\end{table}
